\documentclass[12pt]{article}

\usepackage{geometry}
\geometry{a4paper, margin=1in}

\usepackage{graphicx}
\usepackage{booktabs}
\usepackage{amsmath}
\usepackage{float}
\usepackage{hyperref}

\title{Clinical Data Analysis and Predictive Modeling}
\author{Derek}
\date{\today}

\begin{document}

\maketitle

\section{Introduction}
This report presents a complete data analysis workflow including data loading, preprocessing,
statistical analysis, predictive modeling, and model evaluation using a clinical-related dataset.
The goal is to demonstrate the application of data analysis techniques in a medical context.

\section{Data Source and Description}

The dataset used in this study is the Wisconsin Diagnostic Breast Cancer (WDBC) dataset,
which is a publicly available clinical dataset obtained from the UCI Machine Learning Repository.
The dataset contains measurements computed from digitized images of breast mass biopsies
and is widely used for medical data analysis and disease prediction tasks.

Each observation represents one patient, and the outcome variable indicates whether the tumor
is malignant or benign. The dataset includes 569 patient samples with 30 continuous numerical
features describing tumor characteristics, such as radius, texture, perimeter, and smoothness.

The data were provided in comma-separated values (CSV) format and were loaded into the analysis
environment using the Python \texttt{pandas} library.

\begin{verbatim}
import pandas as pd
data = pd.read_csv("wdbc.csv")
\end{verbatim}

\section{Data Preprocessing}

Data preprocessing was conducted to ensure data quality and suitability for statistical
analysis and predictive modeling. The dataset was first examined for missing values.
No missing data were identified in the selected features, and therefore no imputation
was required.

The diagnosis variable was encoded as a binary outcome, where malignant tumors were
coded as 1 and benign tumors as 0. All continuous features were standardized to have
zero mean and unit variance to improve model convergence and ensure comparability
across variables.

After preprocessing, the dataset was randomly split into training and testing sets,
with 70\% of the data used for model training and 30\% reserved for model evaluation.
The splitting was performed at the patient level to avoid information leakage.

\section{Statistical Analysis}

Descriptive statistical analysis was performed to summarize baseline characteristics
of the study population. Continuous variables were reported as mean and standard deviation.

Comparative statistical analyses were conducted to examine differences between malignant
and benign tumor groups. Independent two-sample t-tests were used for continuous variables
to compare group means. Statistical significance was assessed using a two-sided test
with a significance level of $\alpha = 0.05$.

Several tumor features, including radius, perimeter, and area, showed statistically
significant differences between malignant and benign groups, indicating their potential
value as predictive biomarkers.

\section{Predictive Modeling}

Predictive models were developed to classify tumors as malignant or benign based on
clinical features. Logistic Regression was used as a baseline statistical model due to
its interpretability and suitability for binary classification tasks.

In addition, a Random Forest classifier was implemented to capture potential nonlinear
relationships among features. Model training was performed using the training dataset,
and hyperparameters were selected based on standard practices to balance model performance
and generalization.

\section{Model Evaluation and Visualization}

Model performance was evaluated on the independent testing dataset. Evaluation metrics
included classification accuracy, sensitivity, specificity, and the area under the
receiver operating characteristic curve (ROC-AUC).

ROC curves were used to visualize model discrimination ability, and confusion matrices
were examined to assess classification performance. The Random Forest model demonstrated
improved predictive performance compared to Logistic Regression, indicating its ability
to capture complex feature interactions.

Feature importance analysis further highlighted the contribution of key tumor
characteristics to model predictions.

\section{Discussion}

This study demonstrates that machine learning models can effectively predict breast
cancer malignancy using structured clinical data. Traditional statistical models provide
interpretability, while ensemble methods offer improved predictive performance.

Although this study focused on structured tabular data, medical imaging data such as
mammography or histopathology images may further enhance predictive accuracy when combined
with clinical features.

\section{Conclusion}

This project presents a complete and reproducible clinical data analysis pipeline,
including data preprocessing, statistical analysis, predictive modeling, and model
evaluation. The results highlight the effectiveness of machine learning approaches
in medical decision support.

\end{document}
